(1) 每自旋的二维波矢态密度为$A/(2\pi)^2$，半径$k$内总态数$gAk^2/(4\pi)$。由$\varepsilon=\hbar^2k^2/(2m)$得能量态密度$D=dN(\varepsilon)/d\varepsilon=gAm/(2\pi\hbar^2)$，与能量无关。零温粒子数$N=DE_F$，所以$E_F=N/D$。
(2) Fermi分布$f(\varepsilon)=[e^{(\varepsilon-\mu)/(k_BT)}+1]^{-1}$。粒子数$N=D\int_0^\infty f(\varepsilon)d\varepsilon=Dk_BT\ln(1+e^{\mu/(k_BT)})$。令该式等于$DE_F$后精确解得$\mu(T)=k_BT\ln(e^{E_F/(k_BT)}-1)$，不需要Sommerfeld截断即可得到任意温度结果。
(3) 内能$U=D\int_0^\infty\varepsilon f(\varepsilon)d\varepsilon=-D(k_BT)^2\operatorname{Li}_2[-e^{\mu/(k_BT)}]$。在$T\ll T_F=E_F/k_B$下，Sommerfeld公式给$U=DE_F^2/2+(\pi^2/6)D(k_BT)^2+O(e^{-E_F/(k_BT)})$（固定$N$时指数小修正可乘温度多项式）。故$C_V=(\partial U/\partial T)_{N,A}\simeq(\pi^2/3)Dk_B^2T$，为线性低温热容。
(4) 因$\mu=E_F+k_BT\ln(1-e^{-E_F/(k_BT)})$，二维低温化学势偏离$E_F$的量级是$-k_BT\,e^{-E_F/(k_BT)}$，而非三维自由电子的$-T^2$主项；原因正是二维态密度常数、其能量导数为零。高温$E_F/(k_BT)\ll1$时$\mu\simeq k_BT\ln(E_F/(k_BT))$，对应二维非简并Boltzmann气体。
